% =====================================================================
%  5 -- THE OBSERVABLE END  (source Part 0, section 0.6).
%
%  Source: tier3.md lines 1937-2210, in full and in order --
%          0.6   intro (1937-1941)                    -> \section
%          0.6.1 56Ni and the light curve (1943-1996)  -> sec:ni56
%          0.6.2 Isotopic ratios and the neutron excess (1998-2045)
%                                                      -> sec:isotopicratios
%          0.6.3 Galactic chemical evolution (2047-2068) -> sec:gce
%          0.6.4 Remnants: the two objects everyone cites (2070-2093)
%                                                      -> sec:remnants
%          0.6.5 (warn glyph) The derivative nobody has assembled
%                (2095-2151)                           -> sec:missingderivative
%          0.6.6 The error hierarchy (2153-2208)       -> sec:errorhierarchy
%
%  No self-check block belongs to this file.
%
%  Re-homings: none.  Rewordings: none.  Every internal "SS0.x.y" pointer
%  is rendered as a live \ref through the label registry; Tier 2 pointers
%  ("Tier 2 SS I.5", "Tier 2 SS 0.3.3", "Tier 2 SS 0.3.4", "Tier 2 SS IV.9")
%  stay as plain text; "tier1.md SS VIII.C.1" -> \tierone{VIII.C.1}.
%
%  Source ambiguity resolved: the two "[typical of ...]" tags in 0.6.1 are
%  rendered with \typical[of ...] (the macro prints "[typical; of ...]").
% =====================================================================

\section{The observable end}
\label{part:observables}

The chain has to terminate in something measurable, or the project is
engineering rather than science. This section is what the measurements are,
how they connect to \Ye, and --- at the end --- the derivative nobody has
assembled.

% =====================================================================
\subsection{\texorpdfstring{\nuc{Ni}{56}}{56Ni} and the light curve}
\label{sec:ni56}

The supernova's optical display is powered by radioactive decay of the
iron-peak material the explosion synthesised. The chain:
%
\begin{equation}\label{eq:nichain}
{}^{56}\mathrm{Ni} \xrightarrow{\ t_{1/2} = 6.075\ \mathrm{d}\ } {}^{56}\mathrm{Co}
\xrightarrow{\ t_{1/2} = 77.24\ \mathrm{d}\ } {}^{56}\mathrm{Fe} .
\end{equation}

Both steps are weak, and they are \textbf{not the same weak channel}:
\nuc{Ni}{56} decays by electron capture essentially 100\% (its $\beta^{+}$
branch is $1.3 \times 10^{-5}$), while \nuc{Co}{56} is $\approx$80\% EC and
$\approx$20\% $\beta^{+}$ (ENSDF via the IAEA Live Chart: $\beta^{+}$ 19.7\%,
EC 80.3\%; half-lives 6.075~d and 77.24~d --- ENSDF, retrieved 2026-08-18
\citep{ensdf}). Given how carefully Tier~2 \S I.5 separates the lepton ledgers
for $\beta^{+}$ versus EC, this distinction is worth keeping straight.

The specific energy deposition rates are $\varepsilon_{\mathrm{Ni}} \approx 3.9
\times 10^{10}$ and $\varepsilon_{\mathrm{Co}} \approx 6.8 \times 10^{9}$
erg\,g$^{-1}$\,s$^{-1}$ at $t = 0$ (\citealt{nadyozhin1994}; re-derived here
from the ENSDF mean energies per decay, 1.735 and 3.739~MeV: $3.95 \times
10^{10}$ and $6.70 \times 10^{9}$). \textbf{Arnett's rule} \citep{arnett1982}
--- the statement that at light-curve peak the luminosity approximately equals
the instantaneous radioactive energy deposition rate --- then gives, for
$M(\nuc{Ni}{56}) = 0.07~\Msun$ \typical[of neutrino-driven models above
$\sim$12~\Msun, \citealt{sukhbold2016}; SN 1987A: 0.072--0.077~\Msun], a peak
time $\sim$18 days \typical[of stripped-envelope events; \citealt{lyman2016},
not opened here], and the \emph{mean} lives $\tau = t_{1/2}/\ln 2$
($\tau_{\mathrm{Ni}} = 8.76$~d, $\tau_{\mathrm{Co}} = 111.4$~d;
\citet{diehl2015} quote ``$\tau \sim 8$ days'' and ``$\tau \sim 111$ days''):
%
\begin{equation}\label{eq:lpeak}
L_{\mathrm{peak}} \;\approx\; M_{\mathrm{Ni}}\left[\varepsilon_{\mathrm{Ni}}e^{-t/\tau_{\mathrm{Ni}}} + \varepsilon_{\mathrm{Co}}\left(e^{-t/\tau_{\mathrm{Co}}} - e^{-t/\tau_{\mathrm{Ni}}}\right)\right]
\;\approx\; 1.4\times10^{42}\ \mathrm{erg\ s^{-1}} .
\end{equation}

\derivedhere[$1.38 \times 10^{42}$ recomputed with the mean lives; using the
half-lives in the exponents would give $9.8 \times 10^{41}$; right order of
magnitude for a stripped-envelope event]

\textbf{A qualification that matters and is routinely botched.} Arnett's rule
applies at \emph{peak} to \textbf{stripped-envelope} supernovae (Types Ib/c),
where there is no hydrogen envelope and the light curve is
radioactivity-powered throughout. A 20~\Msun{} progenitor that retains its
envelope produces a \textbf{Type IIP}, whose plateau is powered by hydrogen
recombination \citep{popov1993, kasen2009} --- there \nuc{Ni}{56} sets the
post-plateau radioactive tail and modulates the plateau \emph{length}, not the
peak luminosity \citep{kasen2009, nakar2016}. Tier~2 \S 0.3.3 makes this point;
it is repeated here because ``\nuc{Ni}{56} mass sets the peak brightness'' is a
sentence that will be wrong for most of the events this project's progenitors
produce.

Either way, the connection to \Ye{} is direct: \textbf{\Ye{} sets how much of
the ejecta reaches \nuc{Ni}{56} rather than \nuc{Fe}{54} or \nuc{Ni}{58}}
(\S\ref{sec:endpoint}; \citealt{woosley2007}), so the \nuc{Ni}{56} mass --- a
comparatively well-determined nucleosynthetic quantity (\citealt{sukhbold2016}
call the SN 1987A and Crab \nuc{Ni}{56} masses ``fairly well determined'') ---
is a readout of the neutron excess of the burnt material.

% =====================================================================
\subsection{Isotopic ratios and the neutron excess}
\label{sec:isotopicratios}

The sharper observables are ratios, because they cancel much of the dependence
on the total ejected mass and the mass cut.

\begin{itemize}
\item \textbf{\nuc{Ni}{57}/\nuc{Ni}{56}} (observed as \nuc{Co}{57}/\nuc{Co}{56}
  in the late light curve, since $\nuc{Ni}{57} \to \nuc{Co}{57} \to
  \nuc{Fe}{57}$). \nuc{Ni}{57} has one extra neutron, so its yield relative to
  \nuc{Ni}{56} increases with the neutron excess $\eta = 1 - 2\Ye$
  \assumednote{standard}, but it is nonzero at $\eta = 0$ and set by the
  freeze-out type (\citealt{seitenzahl2014}, Fig.~1, computed at $\Ye = 0.5$
  exactly). Measured in SN 1987A from the light-curve tail shape
  (\citealt{seitenzahl2014}: $2.5 \pm 1.1 \times$ solar) and directly from the
  \nuc{Co}{57} $\gamma$-lines (\citealt{kurfess1992}: $1.5 \pm 0.3 \pm 0.2
  \times$ solar).
\item \textbf{\nuc{Ni}{58}/\nuc{Fe}{56}} and \textbf{\nuc{Fe}{54}/\nuc{Fe}{56}}
  in the solar abundance pattern. The Sun's iron-peak isotopic composition is a
  stringent constraint on the \emph{average} \Ye{} of all the core-collapse
  ejecta that ever enriched the interstellar medium. Overproducing the
  neutron-rich iron-peak isotopes relative to solar is a classic failure mode
  of supernova nucleosynthesis models (\citealt{woosley2007}: a mass cut deeper
  than the iron core gives ``unacceptable overproductions of \nuc{Fe}{54,58}''),
  and it is a direct \Ye{} diagnostic.
\item \textbf{\nuc{Ti}{44}}, produced in $\alpha$-rich freeze-out
  (\S\ref{sec:alpharich}; \citealt{the1998}) and observed as $\gamma$-ray lines
  from Cas~A \citep{grefenstette2014} and SN 1987A \citep{boggs2015}. Sensitive
  to the freeze-out conditions and hence to the shock energetics and mass cut
  as well as to \Ye.
\end{itemize}

The general shape: \textbf{$\eta$ ranges over 0--0.10 across the regime box
(\S\ref{sec:presnye}), and the neutron-rich isotopic yields are roughly
proportional to $\eta$.} The amplification follows from $\eta = 1 - 2\Ye$, so
$\delta\eta = 2\,\delta\Ye$ and
%
\begin{equation}\label{eq:etalever}
\frac{\delta\eta}{\eta} \;=\; \frac{2\,\delta Y_e}{1-2Y_e}
\;=\; \frac{2Y_e}{1-2Y_e}\cdot\frac{\delta Y_e}{Y_e}
\;\approx\; 49\times\frac{\delta Y_e}{Y_e}
\quad\text{at } Y_e = 0.49 .
\end{equation}

So a \emph{fractional} error in \Ye{} of $10^{-3}$ near $\Ye = 0.49$ (i.e.\
$\delta\Ye = 4.9 \times 10^{-4}$) is a \textbf{4.9\% error in $\eta$}, and a
4.9\% error in the yields that depend on it. \emph{That} is the correct way to
see why absolute \Ye{} accuracy at the $10^{-3}$--$10^{-2}$ level is
scientifically meaningful: the small quantity is $\eta$, not \Ye, and near the
neutron-poor end of the box a relative error in \Ye{} is amplified by $\sim$50
into a relative error in the observable. Solar isotopic ratios are known to
percent precision; individual-supernova ratios (\nuc{Ni}{57}/\nuc{Ni}{56} in SN
1987A) to $\sim$20--40\%.

[The $\eta$-linearity claim is an approximation valid for the mildly
neutron-rich isotopes --- and, since the \nuc{Ni}{57}/\nuc{Ni}{56} ratio is
nonzero at $\eta = 0$ \citep{seitenzahl2014}, the $49\times$ lever is an upper
bound on the sensitivity; it is not a substitute for the missing derivative in
\S\ref{sec:missingderivative}. Tagged \textbf{\derivedhere[approximate]}.]

% =====================================================================
\subsection{Galactic chemical evolution}
\label{sec:gce}

Zoom out. Core-collapse supernovae are the dominant source of the $\alpha$
elements (O, Ne, Mg, Si, S, Ca) and a major source of the iron peak; Type Ia
supernovae contribute roughly half of the iron-peak elements, on a delayed
timescale \citep{kobayashi2020}. The observed [$\alpha$/Fe] versus [Fe/H] trend
in Galactic stars --- a plateau at low metallicity followed by a downturn --- is
the fingerprint of that timing difference \citep{tinsley1979, kobayashi2020}.

The relevance here: \textbf{GCE models take supernova yields as inputs,
integrated over an initial mass function} \citep{kobayashi2020}. The yields are
computed from progenitor models whose pre-supernova structure and composition
come from stellar-evolution codes running nuclear networks. A systematic error
in the network's \Ye{} propagates into every yield table, into every GCE model
that uses it, and into the inferred history of the Galaxy. That is a genuinely
large lever, and it is the argument that survives even if one is sceptical
about explodability predictions.

It is also an argument with a specific vulnerability: yield tables are
dominated by uncertainties in the explosion mechanism and mass cut
(\S\ref{sec:masscut}), and the composition improvement this project offers sits
underneath those. The correct claim is \emph{reducing one identified systematic
in the input}, not \emph{fixing yield tables}.

% =====================================================================
\subsection{Remnants: the two objects everyone cites}
\label{sec:remnants}

\textbf{SN 1987A.} The only core-collapse supernova with detected neutrinos
($\sim$24 events: 11 in Kamiokande II, \citealt{hirata1987}; 8 in IMB,
\citealt{bionta1987}; 5 in Baksan, \citealt{alekseev1988}; review
\citealt{janka2017}), which confirmed the basic energetic picture --- total
neutrino energy of a few $\times 10^{53}$~erg (``some $10^{53}$ erg'',
\citealt{janka2017}) over $\sim$10~s (the signal spanned about 12~s), matching
\S\ref{sec:promptshock}'s neutron-star binding energy to within the
uncertainties. Its light curve gave $M(\nuc{Ni}{56}) \approx 0.07~\Msun$
\sourcednote{0.072--0.077~\Msun, Utrobin et al.\ via \citealt{sukhbold2016}},
and its late-time tail constrained \nuc{Ni}{57}/\nuc{Ni}{56}
\citep{seitenzahl2014}.

\textbf{Cassiopeia A.} A $\sim$350-year-old remnant (explosion date $1681 \pm
19$, \citealt{fesen2006}), spatially resolved, with X-ray and $\gamma$-ray maps
of individual elements including \nuc{Ti}{44} \citep{grefenstette2014}. It
shows that the ejecta are strongly asymmetric \citep{milisavljevic2013,
grefenstette2014} and that the iron-peak material is not distributed as a
spherical model would predict --- a reminder that the one-dimensional chain in
\S\ref{sec:corecollapse} is a caricature of a three-dimensional process.

Both are cited in this project's framing as the observational anchors of the
\Ye{} chain. Both deserve the caveat that they constrain the \emph{combination}
of progenitor structure, explosion mechanism and mass cut, and that isolating
the nuclear-network contribution from them has not been done here.

% =====================================================================
\subsection{\texorpdfstring{\warnglyph{}~}{}The derivative nobody has assembled}
\label{sec:missingderivative}

Everything above is qualitative. The quantitative link --- what error in \Ye{}
changes a scientific conclusion --- has never been assembled in this project,
and it is the weakest joint in the motivation.

What exists:
%
\begin{equation}\label{eq:dmchdye-recall}
\frac{dM_{\mathrm{ch}}}{dY_e} \;=\; 5.83\ M_\odot \text{ per unit } Y_e \text{ at } Y_e = 0.5
\qquad\text{(\S\ref{sec:mch}, derived)}
\end{equation}

and the operative per-step gate, $|\Delta\Ye| \lesssim 3 \times 10^{-6}$, which
comes from dividing the \textbf{weak-rate-physics shift} ($5 \times 10^{-3}$ to
$1.5 \times 10^{-2}$ --- \citealt{heger2001}: the amount by which the central
\Ye{} at core collapse of their LMP + $\beta$-decay models exceeds the WW95/FFN
models', about half from adding $\beta$-decays and half from the smaller LMP
capture rates; the project reads it as ``per trajectory'') by $N \approx 1.6
\times 10^{3}$ steps under a systematic-accumulation assumption (Tier~2
\S 0.3.4).

Note what that is: \textbf{an input uncertainty used as an output tolerance.}
The logic --- ``an emulator whose error is below the disagreement between two
defensible rate libraries is not the limiting factor'' --- is sound as a
\emph{floor}, and it is the right way to set one. But it says nothing about
what error would change a conclusion, because
%
\begin{equation}\label{eq:missingderivs}
\frac{\partial M({}^{56}\mathrm{Ni})}{\partial Y_e},\qquad
\frac{\partial\,\xi_{2.5}}{\partial Y_e},\qquad
\frac{\partial(\text{explodability})}{\partial Y_e},\qquad
\frac{\partial(\text{isotopic ratios})}{\partial Y_e}
\end{equation}
%
have not been measured here, cited from the literature, or estimated. The chain
in \S\ref{sec:corecollapse} stops at ``shock birth radius'' and resumes at
``observed \nuc{Ni}{56}'' with nothing connecting them numerically.

Why closing it is worth the day it would take:

\begin{itemize}
\item \textbf{It could move the gate by orders of magnitude in either
  direction.} If explodability is insensitive to $\delta\Ye$ at $10^{-3}$, the
  $3 \times 10^{-6}$ per-step gate is over-engineered by a factor of hundreds
  --- and the project is paying enormously for it (Tier~2 \S IV.9's cost wall
  is downstream of exactly this demand). If some observable is sensitive at
  $10^{-4}$, the gate is right and should be defended loudly.
\item \textbf{It is the referee's first question.} ``You built a $3 \times
  10^{-6}$-per-step emulator. What observable changes at $3 \times 10^{-6}$?''
  The current answer routes through an input uncertainty. That is a defensible
  proxy and it is not an answer.
\item \textbf{Much of it is sourceable rather than computable.} Progenitor and
  explosion model grids in the literature vary \Ye{} (usually implicitly,
  through metallicity, rotation, or rate choices) and report yields. The
  derivative may be extractable from published tables without running anything.
\end{itemize}

Carried as \foreignreg{2}{R-14} and as this tier's \reg{T-01}. It is the
justification for the project's most expensive requirement, and it is missing.

% =====================================================================
\subsection{The error hierarchy --- what the emulator is actually competing with}
\label{sec:errorhierarchy}

\S\ref{sec:missingderivative} asks how accurate the emulator needs to be. This
section asks the complementary and more uncomfortable question: \textbf{how
accurate is everything else?} An improvement that lands two orders below the
dominant error term is real engineering and not a scientific result, and the
only way to know which situation you are in is to assemble the terms.

Ranked by size, for the pre-supernova \Ye{} and composition of a massive-star
model:

\begingroup\footnotesize
\begin{longtable}{@{}L{0.5cm} L{4.5cm} L{5.2cm} L{4.2cm}@{}}
\toprule
\textbf{\#} & \textbf{error term} & \textbf{size (in \Ye{} or its consequences)} & \textbf{status} \\
\midrule
\endfirsthead
\toprule
\textbf{\#} & \textbf{error term} & \textbf{size (in \Ye{} or its consequences)} & \textbf{status} \\
\midrule
\endhead
\midrule
\multicolumn{4}{r}{\emph{continued on next page}}\\
\endfoot
\bottomrule
\endlastfoot
1 & \textbf{3D convection vs MLT} --- entrainment, shell mergers, the
    burning/mixing competition \citep{meakin2007, couch2015, muller2016b}
  & qualitative changes to shell structure and yields; \textbf{not expressible
    as a $\delta\Ye$}
  & \S\ref{sec:convection}; unquantified anywhere \\
2 & \textbf{explosion mechanism + mass cut} (for ejecta quantities only)
  & dominates all ejecta composition claims
  & \S\ref{sec:masscut} \\
3 & \textbf{progenitor parameters} (mass, $Z$, rotation, binarity)
  & the spread across a progenitor grid
  & \S\ref{sec:whichstars} \\
4 & \textbf{weak-rate physics} (WW95/FFN $\to$ LMP + $\beta$-decays)
  & \textbf{$5 \times 10^{-3}$ -- $1.5 \times 10^{-2}$ per model} in central
    \Ye{} at collapse (about half from adding $\beta$-decays, half from smaller
    LMP capture rates)
  & sourced (\citealt{heger2001}), adopted as the project's floor \\
5 & \textbf{network size} (approx21 vs \msn{80} / \msn{151})
  & the term this project removes: the NNN improves \Ye{} over approx21 by
    377--651\% (\msn{80}) / 277--390\% (\msn{151}); the NNN's own per-step
    \Ye{} error is $\approx 2.4 \times 10^{-3}$ absolute (0.33--0.54\% relative
    on \msn{80}), so the approx21 error it improves on is correspondingly
    (4--7$\times$) larger
  & \resultsref{2026-07-08}, \S\ref{sec:nnncritical} \\
6 & \textbf{rate uncertainties within a library} (Hauser--Feshbach factor
    $\sim$2 on unmeasured channels; \citealt{rauscher2000}: ``within a factor
    of 1.5--2'')
  & unpropagated
  & \tierone{VIII.C.1} \\
7 & \textbf{emulator target}
  & $3 \times 10^{-6}$ per step $\to$ $5 \times 10^{-3}$ per trajectory
  & the gate \\
8 & \textbf{conservation drift}
  & $\le 10^{-12}$ per step
  & measured 0.0 \\
\end{longtable}
\endgroup

Read the table honestly and three things follow.

\textbf{(i) The project's gate is set \emph{equal to} term 4 and is therefore
correctly placed relative to the input physics.} The design logic --- do not be
the limiting factor --- is sound, and rows 7 and 4 agreeing at the trajectory
level is exactly the intended relationship. Row 8 being six orders below row 7
is why conservation is free.

\textbf{(ii) But rows 1--3 are larger and are \emph{not} reduced by anything
this project does.} The defensible claim is therefore narrow and should be
stated in this form: \textbf{the project removes term 5 and pushes term 7 below
term 4.} It does not improve the pre-supernova model's total accuracy, because
that is dominated by terms 1--3. What it \emph{does} is remove a term that is
(a) currently large, (b) entirely artificial --- a compute compromise rather
than a physical uncertainty --- and (c) removable at no accuracy cost, which
none of terms 1--3 are.

That distinction --- \textbf{artificial versus physical error} --- is the
strongest framing available and it is not the one the project currently uses.
An artificial error is one you would not make if compute were free. Removing
artificial error is exactly what surrogate modelling is \emph{for}, and it does
not require beating the physical uncertainties to be worth doing.

\textbf{(iii) Term 6 is the one that could embarrass the project.} If
unmeasured Hauser--Feshbach rates carry factor-2 uncertainties on channels that
matter for \Ye, then a $3 \times 10^{-6}$ emulator and a $3 \times 10^{-4}$
emulator produce indistinguishable science, and the gate is over-engineered by
two orders. \tierone{VIII.C.1} has been open since Tier~1 and the machinery to
close it (perturb rates, re-integrate, measure $\delta\Ye$) is entirely built
--- the batch engine evaluates every rate and the reference integrator exists.

\textbf{Term 6 is the cheapest of the three ways to answer ``is the gate
right?''} --- the other two being \reg{T-01} (the output derivative, partly
infeasible per \S\ref{sec:explodability}) and \foreignreg{2}{R-11} (the
operator-split comparison). It should be first.

Registered as \reg{T-20} (assemble the hierarchy as a project document) and it
promotes \tierone{VIII.C.1} to the top of the cross-tier list.
